% Options for packages loaded elsewhere
\PassOptionsToPackage{unicode}{hyperref}
\PassOptionsToPackage{hyphens}{url}
\documentclass[
]{article}
\usepackage{xcolor}
\usepackage[margin=1in]{geometry}
\usepackage{amsmath,amssymb}
\setcounter{secnumdepth}{-\maxdimen} % remove section numbering
\usepackage{iftex}
\ifPDFTeX
  \usepackage[T1]{fontenc}
  \usepackage[utf8]{inputenc}
  \usepackage{textcomp} % provide euro and other symbols
\else % if luatex or xetex
  \usepackage{unicode-math} % this also loads fontspec
  \defaultfontfeatures{Scale=MatchLowercase}
  \defaultfontfeatures[\rmfamily]{Ligatures=TeX,Scale=1}
\fi
\usepackage{lmodern}
\ifPDFTeX\else
  % xetex/luatex font selection
\fi
% Use upquote if available, for straight quotes in verbatim environments
\IfFileExists{upquote.sty}{\usepackage{upquote}}{}
\IfFileExists{microtype.sty}{% use microtype if available
  \usepackage[]{microtype}
  \UseMicrotypeSet[protrusion]{basicmath} % disable protrusion for tt fonts
}{}
\makeatletter
\@ifundefined{KOMAClassName}{% if non-KOMA class
  \IfFileExists{parskip.sty}{%
    \usepackage{parskip}
  }{% else
    \setlength{\parindent}{0pt}
    \setlength{\parskip}{6pt plus 2pt minus 1pt}}
}{% if KOMA class
  \KOMAoptions{parskip=half}}
\makeatother
\usepackage{graphicx}
\makeatletter
\newsavebox\pandoc@box
\newcommand*\pandocbounded[1]{% scales image to fit in text height/width
  \sbox\pandoc@box{#1}%
  \Gscale@div\@tempa{\textheight}{\dimexpr\ht\pandoc@box+\dp\pandoc@box\relax}%
  \Gscale@div\@tempb{\linewidth}{\wd\pandoc@box}%
  \ifdim\@tempb\p@<\@tempa\p@\let\@tempa\@tempb\fi% select the smaller of both
  \ifdim\@tempa\p@<\p@\scalebox{\@tempa}{\usebox\pandoc@box}%
  \else\usebox{\pandoc@box}%
  \fi%
}
% Set default figure placement to htbp
\def\fps@figure{htbp}
\makeatother
\setlength{\emergencystretch}{3em} % prevent overfull lines
\providecommand{\tightlist}{%
  \setlength{\itemsep}{0pt}\setlength{\parskip}{0pt}}
\usepackage{bookmark}
\IfFileExists{xurl.sty}{\usepackage{xurl}}{} % add URL line breaks if available
\urlstyle{same}
% fallback for those not using the hyperref driver hyperxmp:
\makeatletter
\@ifundefined{xmpquote}{\newcommand{\xmpquote}[1]{#1}}{}
\makeatother
\hypersetup{
  pdftitle={PSTAT 132 - Homework 2},
  pdfauthor={Aldo Nunez},
  hidelinks,
  pdfcreator={LaTeX via pandoc}}

\title{PSTAT 132 - Homework 2}
\author{Aldo Nunez}
\date{2026-10-14}

\begin{document}
\maketitle

\textbf{Question 1}

Relational algebra: \(\\\)
\(\pi_{\text{CourseID}}\left(\sigma_{\text{StudentID}='1001'}(\text{ENROLLMENT})\right) - \pi_{\text{CourseID}}\left(\sigma_{\text{StudentID}='1002'}(\text{ENROLLMENT})\right)\)

SQL:

SELECT CourseID FROM ENROLLMENT WHERE STUDENT ID = `1001' \(\\\) EXCEPT
\(\\\) SELECT CourseID FROM ENROLLMENT WHERE STUDENT ID = `1002'

\textbf{Question 2}

Relational algebra:

\(\pi_{\text{Name},\,\text{Title}}\left(\sigma_{\text{Grade}='A'}\left(\text{STUDENT} \bowtie \text{ENROLLMENT} \bowtie \text{COURSE}\right)\right)\)

SQL:

SELECT S.Name, C.Title \(\\\) FROM STUDENT S \(\\\) JOIN ENROLLMENT E ON
S.StudentID = E.StudentID \(\\\) JOIN COURSE C ON E.CourseID =
C.COURSEID \(\\\) WHERE E.Grade = `A'; \(\\\)

\textbf{Question 3}

Relational algebra: \(\\\)
\(\pi_{\text{Name},\,\text{Title}}\left(\text{STUDENT} \bowtie \text{ENROLLMENT} \bowtie \text{COURSE}\right)\)

SQL:

SELECT S.Name, C.Title \(\\\) FROM STUDENT S \(\\\) JOIN ENROLLMENT E ON
S.StudentID = E.StudentID \(\\\) JOIN COURSE C ON E.CourseID =
C.COURSEID \(\\\)

\textbf{Question 5}

sql\_query\_5 \textless- ``WITH customer\_totals AS ( SELECT
customer\_id, SUM(order\_amount) AS total\_spend FROM orders WHERE
order\_date \textgreater= `2023-01-01' AND order\_date \textless{}
`2023-02-01' GROUP BY customer\_id ), avg\_spend AS ( SELECT
AVG(total\_spend) AS avg\_total\_spend FROM customer\_totals ) SELECT
ct.customer\_id, ct.total\_spend, a.avg\_total\_spend, ct.total\_spend -
a.avg\_total\_spend AS diff\_from\_avg, CASE WHEN ct.total\_spend
\textgreater{} a.avg\_total\_spend THEN `Above average' WHEN
ct.total\_spend \textless{} a.avg\_total\_spend THEN `Below average'
ELSE `Equal to average' END AS comparison FROM customer\_totals ct CROSS
JOIN avg\_spend a ORDER BY ct.total\_spend DESC;''

\textbf{Question 6} sql\_query\_6 \textless- ``CREATE VIEW
customer\_monthly\_spend AS SELECT customer\_id, strftime(`\%Y-\%m',
order\_date) AS order\_month, SUM(order\_amount) AS total\_spend FROM
orders GROUP BY customer\_id, strftime(`\%Y-\%m', order\_date);''

\textbf{Question 7} Two intermediate relations: (a) each customer's
total spend for January 2023 (filter by date, GROUP BY customer\_id,
SUM(order\_amount)), and (b) the average of those totals.

The CTE names the per-customer January totals once at the top of the
query. The main query can then read from it twice: once to list
customers and once to compute the average to compare against. Each
logical step is separate and readable. Repeating the aggregation in two
nested subqueries duplicates logic. That is harder to read, and if one
copy is edited (e.g., the date range), the two copies can silently drift
apart. The CTE is a single source of truth. No.~A CTE is temporary and
exists only while that one query executes. Nothing is saved, which is
why SELECT * FROM january\_completed\_orders fails in Demo 3.

No.~A view adds a permanent schema object, which pays off when logic is
reused across many queries or sessions. For a one-off or occasional
analysis, a CTE is simpler and leaves nothing behind to maintain.

\textbf{Question 8}

sql\_query\_8 \textless- ``WITH student\_avg AS ( SELECT StudentID,
AVG(Grade) AS average\_grade FROM ENROLLMENT GROUP BY StudentID HAVING
AVG(Grade) \textgreater= 85 ) SELECT s.StudentID, s.Name, s.Major,
sa.average\_grade FROM student\_avg sa JOIN STUDENT s ON s.StudentID =
sa.StudentID ORDER BY sa.average\_grade DESC;''

The CTE separates the aggregation and filtering step (average grade,
keep only ≥ 85) from the join to STUDENT. The query reads in the same
order as the problem's steps. The join also only touches the
already-reduced set of successful students instead of mixing GROUP
BY/HAVING with the join logic.

\end{document}
